\section{D'où viennent les courbes aberrantes ?}\label{sec-diagnostic}

\begin{objectif}
(O1) Décomposer le repo implicite, calculer la taille du bruit normal, et donner la forme que laisse dans l'écart de
portage chacun des mécanismes qui produisent de vraies anomalies.
\end{objectif}

\subsection{L'écart de panel}\label{sec-gap}
On se place sous les \cref{hyp-coherent,hyp-panels,hyp-consensus} et à une maturité $T$ fixée. Les vues des
contributeurs sont des perturbations d'un centre commun, $q_j=q+u_j$ et $D_j=D+v_j$, où $u_j$ et $v_j$ sont d'ordre
$\epsilon$. On note $\qbar$ et $\Dopt$ les moyennes des $q_j$ et des $D_j$ sur le panel options, et $\Ddiv$ le
consensus du panel dividendes, alloué à la maturité $T$.

\begin{proposition}[Décomposition du repo implicite]\label{prop-gap}
\begin{equation}\label{eq-gap}
q^{\tot}(T)=\qbar(T)+\frac{\Dopt(T)-\Ddiv(T)}{T\,\bigl(S_0-D(T)\bigr)}+O(\epsilon^2).
\end{equation}
\end{proposition}
\begin{proof}
Pour $j\in\Popt$, le \cref{lem-logf} donne $\ln F_j=\ln F-T\,u_j-v_j/(S_0-D)+O(\epsilon^2)$. Le forward Totem est la
moyenne des $F_j$, et le logarithme d'une moyenne est égal à la moyenne des logarithmes à un terme d'ordre $\epsilon^2$
près ; donc $\ln F^{\tot}=\ln F-T\,\bar u^{\opt}-\bar v^{\opt}/(S_0-D)+O(\epsilon^2)$. Le contrôleur retire les
dividendes de consensus $D^{\tot}=D+\bar v^{\dv}$ :
\[
q^{\tot}T=rT-\ln F^{\tot}+\ln\bigl(S_0-D^{\tot}\bigr)=qT+T\,\bar u^{\opt}+\frac{\bar v^{\opt}-\bar v^{\dv}}{S_0-D}
+O(\epsilon^2).
\]
On divise par $T$, avec $q+\bar u^{\opt}=\qbar$ et $\bar v^{\opt}-\bar v^{\dv}=\Dopt-\Ddiv$.
\end{proof}

Le repo implicite est donc le repo moyen du panel options, plus un \emph{écart de panel} : la différence entre la vue
moyenne des dividendes du panel options, que personne n'observe, et celle du panel dividendes, divisée par
$T(S_0-D)$. Pour un contributeur donné, une erreur sur ses dividendes se compense exactement dans son repo implicite ;
entre deux panels, elle ne se compense pas.

\begin{corollary}[Taille de l'écart de panel]\label{cor-gap}
On suppose les erreurs de vue sur $D(T)$ indépendantes, de même loi et d'écart-type $\sigma_D(T)$ dans les deux
panels, et les $n^{\com}$ contributeurs communs aux deux services cohérents avec eux-mêmes. Alors
\begin{equation}\label{eq-gap-sd}
\sigma\bigl(q^{\tot}-\qbar\bigr)=\frac{\sigma_D(T)}{T\,\bigl(S_0-D(T)\bigr)}
\sqrt{\frac{1}{n^{\opt}}+\frac{1}{n^{\dv}}-\frac{2\,n^{\com}}{n^{\opt}n^{\dv}}} .
\end{equation}
\end{corollary}
\begin{proof}
$\Dopt-\Ddiv=\sum_j(\alpha_j-\gamma_j)\,v_j$, avec $\alpha_j=1/n^{\opt}$ sur le panel options, $\gamma_j=1/n^{\dv}$ sur
le panel dividendes, et $0$ ailleurs. Sa variance vaut
$\sigma_D^2\sum_j(\alpha_j-\gamma_j)^2=\sigma_D^2\bigl[n^{\com}(1/n^{\opt}-1/n^{\dv})^2+(n^{\opt}-n^{\com})/(n^{\opt})^2
+(n^{\dv}-n^{\com})/(n^{\dv})^2\bigr]$, qui se simplifie en
$\sigma_D^2\bigl(1/n^{\opt}+1/n^{\dv}-2n^{\com}/(n^{\opt}n^{\dv})\bigr)$.
\end{proof}

L'écart disparaît si les deux panels sont identiques ($n^{\com}=n^{\opt}=n^{\dv}$) ; il est maximal s'ils sont
disjoints. La dispersion $\sigma_D(T)$ d'un dividende cumulé croît au moins comme $T$ : l'écart de panel est petit en
portage sur le court terme, mais il ne diminue pas en repo sur le long terme.

\begin{example}[Taille de l'écart pour une action]\label{ex-gap}
On reprend le modèle de l'exemple fil rouge : spot $100$, dividende annuel de $\num{4.2}$ détaché début mai, et chaque
contributeur se trompe sur le dividende détaché en $t$ d'une erreur relative d'écart-type $\num{0.01}+\num{0.03}\,t$,
dont une part est commune à tous ses dividendes. Les panels sont disjoints, de la taille de ceux de l'exemple fil
rouge, et les contributeurs options divergent aussi de $15$\,pb sur le repo, ce que le repo implicite ne doit pas
capter. La formule \eqref{eq-gap-sd} et une simulation contributeur par contributeur donnent :
\begin{center}\input{generated/\lang/tab_gap}\end{center}
L'écart seul produit \vGapOneY\,pb de bruit sur le repo à un an et \vGapFiveY\,pb à cinq ans, sans aucune erreur dans
Totem. Il est nul avant le premier détachement, saute à chaque détachement, puis décroît en $1/T$ jusqu'au suivant
(\cref{fig-gap}).
\end{example}

\begin{figure}[ht]
\centering\includegraphics[width=\linewidth]{figures/\lang/panel_gap.pdf}
\caption{(a) Bruit du repo implicite dû au seul écart de panel, formule \eqref{eq-gap-sd}, pour les panels de
l'exemple fil rouge et deux autres tailles (options / dividendes) : il naît à chaque détachement puis décroît en $1/T$.
(b) Un décalage $\varepsilon$ du spot est constant en portage, donc vaut $\varepsilon/T$ en repo
(\cref{prop-signatures}).}
\label{fig-gap}
\end{figure}

\subsection{Les sorties de contributeurs}\label{sec-nested}
Les panels se réduisent avec la maturité. Comme ils sont emboîtés, la moyenne sur un petit panel n'est pas
indépendante de la moyenne sur un grand panel qui le contient.

\begin{proposition}[Moyennes sur des panels emboîtés]\label{prop-nested}
Soit $e_{jk}$ l'erreur du contributeur $j$ à la maturité $T_k$, indépendante d'un contributeur à l'autre, avec
$\Var e_{jk}=\sigma_k^2$ et $\Cov(e_{jk}, e_{jl})=\rho_{kl}\,\sigma_k\sigma_l$. Si les panels sont emboîtés, de tailles
$n_k$, les moyennes $\bar e_k$ vérifient
\begin{equation}\label{eq-nested}
\Cov(\bar e_k, \bar e_l)=\frac{\rho_{kl}\,\sigma_k\sigma_l}{\max(n_k, n_l)} .
\end{equation}
En particulier, pour des vues persistantes ($\rho_{kl}=1$, $\sigma_k=\sigma_l=\sigma$) et $n_l\le n_k$, le saut de la
moyenne quand le panel passe de $n_k$ à $n_l$ contributeurs a pour écart-type $\sigma\sqrt{1/n_l-1/n_k}$.
\end{proposition}
\begin{proof}
Les deux moyennes ont $\min(n_k, n_l)$ contributeurs en commun, donc
\[
\Cov(\bar e_k, \bar e_l)=\frac{\min(n_k, n_l)}{n_kn_l}\,\rho_{kl}\sigma_k\sigma_l
=\frac{\rho_{kl}\,\sigma_k\sigma_l}{\max(n_k, n_l)} .
\]
Ensuite $\Var(\bar e_l-\bar e_k)=\sigma^2(1/n_l+1/n_k-2/n_k)=\sigma^2(1/n_l-1/n_k)$.
\end{proof}

En unités de dispersion des contributeurs, le consensus saute ainsi de lui-même de $\vJumpA\,\sigma$ quand le panel
passe de $16$ à $12$ contributeurs, de $\vJumpB\,\sigma$ de $12$ à $7$ (le panel options de l'exemple fil rouge entre
deux et trois ans), de $\vJumpC\,\sigma$ de $15$ à $7$ et de $\vJumpD\,\sigma$ de $10$ à $4$. Un saut de cette taille à
une sortie de contributeurs est du bruit, pas une information ; le modèle de bruit de la \cref{sec-detection} le
reproduit grâce à \eqref{eq-nested}.

\subsection{Six mécanismes, six signatures}\label{sec-signatures}
Au-delà du bruit, certains mécanismes produisent de vraies anomalies. Chacun laisse dans l'écart de portage une forme
caractéristique.

\begin{proposition}[Signatures]\label{prop-signatures}
Au premier ordre, chaque mécanisme ajoute à l'écart de portage $\delta$ la contribution suivante.
\begin{enumerate}[label=(\roman*)]
\item \emph{Décalage du spot.} Si les contributeurs options valorisent avec un spot $S_0(1+\varepsilon)$ alors que le
contrôleur utilise $S_0$ : $-\varepsilon S_0/(S_0-D_k)\approx-\varepsilon$ à tous les piliers, soit une erreur de
repo $-\varepsilon/T_k$. Un dividende détaché avant $T_1$ et erroné de $\delta d$ dans le service dividendes produit la
même forme, $-\delta\tilde d/(S_0-D_k)$.
\item \emph{Erreur sur une période de dividendes.} Si le niveau Totem de la période $p$ est erroné de $\delta\ell_p$ :
$-a_{kp}\,\delta\ell_p/(S_0-D_k)$, une marche qui commence au premier pilier qui suit le détachement ; en repo, un saut
qui décroît comme $1/T$.
\item \emph{Désaccord sur la date de détachement.} Soit un dividende détaché en $t_i\in\intoo{T_k}{T_{k+1}}$, et une
fraction $f$ du panel options qui le place avant $T_k$ : $+f\,\tilde d_i/(S_0-D_k)$ au seul pilier $T_k$, un pic. La
variance des $\ln F_j$ en $T_k$ augmente de $f(1-f)\,\tilde d_i^{\,2}/(S_0-D_k)^2$.
\item \emph{Profil de détachement périmé.} Si les forwards utilisent les vraies dates $t_i$ mais l'allocation les
place en $t_i+\Delta$ avec $\Delta>0$ : $+\tilde d_i/(S_0-D_k)$ aux piliers $T_k\in\intco{t_i}{t_i+\Delta}$, et rien
ailleurs (pour $\Delta<0$, $-\tilde d_i/(S_0-D_k)$ aux piliers $T_k\in\intco{t_i+\Delta}{t_i}$). Pour un payeur
saisonnier, les piliers touchés reviennent chaque année.
\item \emph{Forward aberrant.} Si le consensus $\ln F^{\tot}_k$ est décalé de $\mu$ : $-\mu$ au seul pilier $T_k$,
un pic.
\item \emph{Biais proportionnel sur les dividendes} (par exemple brut contre net de retenue à la source). Si chaque
dividende Totem vaut $(1+\tau)$ fois celui des forwards : $-\tau D_k/(S_0-D_k)$, un escalier qui monte à chaque
détachement. Pour un flux continu de rendement $y$, $D_k\approx S_0\,y\,T_k$ et la contribution vaut environ
$-\tau\,y\,T_k$ : exactement la forme d'un écart de repo constant de $-\tau y$.
\end{enumerate}
Deux autres formes complètent le tableau : une sortie de contributeurs au pilier $T_k$ produit une marche en $T_k$
(\cref{prop-nested}), et une erreur de marquage du repo FO de $\Delta q$ produit $-\Delta q(T_k)\,T_k$, un écart de
repo lisse.
\end{proposition}
\begin{proof}
On applique \eqref{eq-delta-exact} à forwards ou dividendes perturbés.
(i) Le forward de chaque contributeur devient $\bigl(S_0(1+\varepsilon)-D\bigr)e^{(r-q)T}$, donc $\ln F^{\tot}_k$ varie
de $\ln\bigl(1+\varepsilon S_0/(S_0-D_k)\bigr)$ ; pour le dividende proche, $D^{\tot}_k$ varie de $\delta\tilde d$ à
tous les piliers.
(ii) Le \cref{lem-allocation} donne $\delta D^{\tot}_k=a_{kp}\,\delta\ell_p$, nul avant le détachement et à peu près
constant après.
(iii) En $T_k$, une fraction $f$ du panel options inclut le dividende, donc $\ln F^{\tot}_k$ baisse de
$f\tilde d_i/(S_0-D_k)$, alors que l'allocation l'exclut ; avant $T_k$, personne ne l'inclut, et après $T_{k+1}$ tout
le monde l'inclut. Pour la variance : les contributeurs forment deux groupes, de poids $f$ et $1-f$, dont les $\ln F_j$
moyens diffèrent de $\tilde d_i/(S_0-D_k)$.
(iv) $D^{\tot}_k-D_k=\sum_i\tilde d_i\bigl(\one_{\{t_i+\Delta\le T_k\}}-\one_{\{t_i\le T_k\}}\bigr)$, qui vaut
$-\tilde d_i$ exactement quand $t_i\le T_k<t_i+\Delta$.
(v) Immédiat. (vi) $\delta D^{\tot}_k=\tau D_k$ ; pour un rendement, $D_k=S_0(1-e^{-yT_k})\approx S_0\,y\,T_k$ et
$S_0-D_k\approx S_0$.
\end{proof}

\begin{table}[ht]
\centering
\caption{Signatures des mécanismes. La dernière colonne donne la forme type que teste le détecteur de la
\cref{sec-detection}.}
\label{tab-signatures}
\small
\begin{tabular}{@{}>{\raggedright\arraybackslash}p{4.2cm}>{\raggedright\arraybackslash}p{4.3cm}
>{\raggedright\arraybackslash}p{3.6cm}l@{}}
\toprule
Mécanisme & Écart de portage $\delta$ & Écart de repo & Forme type\\
\midrule
Écart de panel (\ref{prop-gap}) & bruit lisse, incréments aux détachements & bruit à peu près plat & bruit\\
Sortie de contributeurs (\ref{prop-nested}) & marche à la sortie & saut décroissant en $1/T$ & marche\\
(i) Décalage du spot & constante & hyperbole $-\varepsilon/T$ & niveau\\
(ii) Erreur sur une période & marche après le détachement & saut décroissant en $1/T$ & marche\\
(iii) Désaccord de date & un pilier & pic & pic\\
(iv) Profil périmé & piliers proches des détachements & zigzag annuel & saison\\
(v) Forward aberrant & un pilier & pic & pic\\
(vi) Biais proportionnel & escalier ; $\approx-\tau yT$ pour un flux & marches ; niveau pour un indice & marche\\
Erreur de marquage FO & $-\Delta q(T)\,T$, lisse & niveau ou pente lisse & aucune\\
\bottomrule
\end{tabular}
\end{table}

\begin{filrouge}
La \cref{fig-gallery} montre l'action de l'exemple fil rouge sous chacun des huit scénarios du banc d'essai, avec le
même bruit de contributeurs que dans le marché sain. Le spot des contributeurs options est décalé de \vRunSpotEps{} en
(a) : la courbe Totem plonge en $1/T$ sur le court terme. En (b), le niveau Totem 2027 est décalé de
\vRunBucketDelta{} : saut à neuf mois, qui s'amortit en $1/T$. En (c), \vRunTimingShare{} du panel options placent avant
l'échéance de six mois le dividende détaché \vRunTimingDays~jours après : pic à six mois. En (d), le profil d'allocation
est décalé de \vRunSeasonShift~jours : les dividendes de mai passent avant les échéances de six et dix-huit mois, d'où
deux pics négatifs. En (e), le consensus $\ln F$ à deux ans est décalé de \vRunFwdShift\,pb. En (f), les contributeurs
options restants après la sortie de trois ans surestiment de \vRunPanelDelta{} les dividendes suivants : la
déformation est lisse et reste dans la bande. En (g), c'est le repo FO qui est décalé de \vRunFoShift\,pb, Totem étant
intact. En (h), toutes les périodes Totem sont biaisées de \vRunGrossBias{} (biais brut/net).
\end{filrouge}

\begin{figure}[p]
\centering\includegraphics[width=\linewidth]{figures/\lang/gallery.pdf}
\caption{Exemple fil rouge sous les huit scénarios : repo FO, repo implicite Totem, bande de $\pm2$ écarts-types du
bruit de consensus autour du FO, et, en vert, la courbe Totem réparée par le détecteur de la \cref{sec-detection}
(anomalie retirée). En (f), le détecteur ne signale rien ; en (g), il ne signale pas d'anomalie Totem mais une erreur de
marquage FO possible. Les diagnostics sont détaillés au \cref{tab-running}.}
\label{fig-gallery}
\end{figure}

\subsection{Ce que les données ne peuvent pas dire}\label{sec-nonid}
\begin{theorem}[Non-identifiabilité de la répartition]\label{thm-nonid}
On suppose que les prix d'options européennes ne dépendent des vues d'un contributeur qu'à travers ses forwards.
Pour tout contributeur $j$ du panel options, on remplace $(q_j, D_j)$ par $(q_j+u_j, D'_j)$ avec
$u_j(T)\,T=\ln\bigl(S_0-D'_j(T)\bigr)-\ln\bigl(S_0-D_j(T)\bigr)$. Alors tous les prix d'options sont inchangés, donc
toutes les statistiques du panel options (consensus, dispersion, nombre de contributeurs) ; les statistiques du panel
dividendes le sont aussi. Le repo implicite $q^{\tot}$ est donc inchangé, alors que le repo moyen du panel options
$\qbar$ varie de la moyenne des $u_j$. Par conséquent, l'écart de panel $q^{\tot}-\qbar$ n'est pas une fonction des
données observables.
\end{theorem}
\begin{proof}
La transformation laisse $\ln F_j=(r-q_j)T+\ln(S_0-D_j)$ inchangé pour tout $j$ et tout $T$ : les prix d'options de
chaque contributeur, et toute statistique qui en découle, sont inchangés. Le panel dividendes n'est pas touché.
$q^{\tot}$ n'est calculé qu'à partir de ces statistiques ; il ne bouge donc pas, alors que $\qbar$ se déplace de la
moyenne des $u_j$, qui peut prendre n'importe quelle valeur.
\end{proof}

\begin{corollary}\label{cor-nonid}
Aucun estimateur construit à partir des seules données de consensus Totem (lissage de courbe, regroupement, réseau
de neurones) ne peut retirer l'écart de panel. On peut en revanche en calculer la taille (\cref{sec-detection}),
reconnaître les mécanismes structurés à leur signature (\cref{prop-signatures}), et se passer de la répartition
entre repo et dividendes là où l'exposition ne l'exige pas (\cref{sec-controle}).
\end{corollary}

\begin{remark}[Ce qui lève l'indétermination]
Des informations extérieures au consensus fixent directement l'un des deux côtés : contrats à terme cotés sur
dividendes, contrats à terme de rendement total, qui cotent un spread de financement, commissions de prêt-emprunt de
titres pour le court terme. Des contributeurs communs aux deux services réduisent l'écart (\cref{cor-gap}), et les
options américaines portent un peu d'information sur la répartition par l'exercice anticipé.
\end{remark}

\begin{retenir}
\begin{itemize}[leftmargin=12pt]
\item Repo implicite $=$ repo du panel options $+$ écart de panel. L'écart de panel est du bruit normal, de taille
calculable : quelques points de base à un an, une à plusieurs dizaines à cinq ans.
\item Les vraies anomalies ont des formes reconnaissables en portage : niveau (spot), marche (période de dividendes,
sortie de contributeurs), pic (date de détachement, forward aberrant), saison (profil périmé).
\item La répartition entre repo et dividendes du panel options n'est pas identifiable à partir de Totem : aucune
méthode statistique, si sophistiquée soit-elle, ne la retrouvera.
\end{itemize}
\end{retenir}
